\documentclass[10pt,a4paper]{amsart}
\usepackage[margin=19mm]{geometry}
\usepackage{amsmath,amssymb,mathtools}
\usepackage[scr=rsfs,cal=euler]{mathalfa}
\usepackage{xfrac}
\usepackage[unicode,hidelinks,bookmarks=false]{hyperref}
\newcommand{\ie}{i.e.\ }
\newcommand{\cf}{cf.\ }
\newcommand{\resp}{respectively\ }
\newcommand{\Subsection}{\subsection}
\providecommand{\nobreakdash}{\nobreak\hbox{-}}
\setlength{\emergencystretch}{2em}
\multlinegap0pt
\usepackage{amssymb}
\usepackage{float}
\usepackage{bm}
\newtheorem{cor}{Corollary}
\title{A new family of Gaussian processes for modeling animal movement: application to bat telemetry data}
\author{J.H. Ram\'{\i}rez-Gonz\'{a}lez, A. Murillo-Salas, Ying Sun}
\date{Source: arXiv:2405.19903v4 (CC BY 4.0)}
\begin{document}
\maketitle

\noindent\emph{Source and licence.} A new family of Gaussian processes for modeling animal movement: application to bat telemetry data. Version arXiv:2405.19903v4 (CC BY 4.0), distributed by arXiv under the Creative Commons Attribution 4.0 International licence, http://creativecommons.org/licenses/by/4.0/, as recorded in the arXiv licence metadata for this exact version and shown on the abstract page https://arxiv.org/abs/2405.19903v4. Full licence notice: \texttt{licenses/arxiv-2405.19903-CC-BY-4.0.txt}.\par\smallskip

\noindent\textit{Editorial scope.} Counted unit: the article's cor cor:stein-power-closure (source0.tex lines 3775--3980). The article's complete original proof of it is delivered with it (source0.tex lines 3982--4167, one \texttt{proof} environment of 186 lines). No mathematics is written, repaired or re-derived: the statement and the proof are the article's own lines, in the article's own notation, and no retained line is substituted or re-flowed.  Retained uncounted as the source-local context the proof needs: lines 3674--3680.  Nothing was omitted from any retained span, so the package carries no omission record.  No retained line contains a citation, so the wrapper carries no bibliography and no bibliography entry.  No retained line contains a figure environment, an image-inclusion command or drawing code, so no image is delivered and the package declares no asset dependency.  Editorial formatting only, in addition: an AMS \texttt{amsart} preamble that restates the article's own declarations and macros used by the retained lines, namely the environments \texttt{cor} on separate counters, together with the standard packages those lines need.  The retained environments print in this document as Corollary 1 in order of appearance, whereas the article prints the same items with the numbering of its own full document.  The complete native source of the article as distributed in the arXiv e-print for this exact version is bundled unchanged as \texttt{sources/160008/source0.tex}.
\begin{equation}\label{eq:increment-double-integral}
D_T
=
-\int_r^v\int_s^t
R_{d,\nu,\rho}''(T+y-x)
\,dy\,dx.
\end{equation}

\begin{cor}
\label{cor:stein-power-closure}
Let \(d\in\mathbb N\), with \(\mathbb N=\{1,2,\ldots\}\), and set
\[
p:=\frac d2.
\]
For \(\nu>0\), \(\rho>0\), and \(\sigma^2>0\), consider the
LRD-compatible Stein temporal covariance
\[
R_{d,\nu,\rho}(h)
:=
\sigma^2
\frac{\Gamma(\nu+p)}{\Gamma(\nu)}
\frac{\Gamma(\nu+\rho|h|)}
{\Gamma(\nu+\rho|h|+p)},
\qquad h\in\mathbb R.
\]
For \(\lambda>0\), define
\begin{equation}\label{eq:Sd-covariance}
R^{S,d}_{\sigma^2,\lambda}(h)
:=
\frac{\sigma^2}{(1+\lambda|h|)^p},
\qquad h\in\mathbb R.
\end{equation}
Then the following statements hold.

\begin{enumerate}
\item[(i)]
The function \(R^{S,d}_{\sigma^2,\lambda}\) is a valid stationary
covariance function on \(\mathbb R\). Accordingly, let
\(S^{(d)}_{\sigma^2,\lambda}\) denote a centered stationary Gaussian
process with covariance \eqref{eq:Sd-covariance}.

\item[(ii)]
The process \(S^{(d)}_{\sigma^2,\lambda}\) exhibits LRD. More precisely,
for every \(r<v\) and \(s<t\),
\begin{equation}\label{eq:Sd-lrd-limit}
\begin{aligned}
&\lim_{T\to\infty}T^{p+2}
\operatorname{Cov}\left(
S^{(d)}_{\sigma^2,\lambda}(v)
-
S^{(d)}_{\sigma^2,\lambda}(r),
\right.\\[-1mm]
&\hspace{4.3cm}\left.
S^{(d)}_{\sigma^2,\lambda}(t+T)
-
S^{(d)}_{\sigma^2,\lambda}(s+T)
\right)\\
&\qquad
=
-\sigma^2\lambda^{-p}p(p+1)(v-r)(t-s).
\end{aligned}
\end{equation}

\item[(iii)]
Under the parametrization
\[
\rho=\nu\lambda,
\]
the Stein family has two nondegenerate boundary closures. As
\(\nu\to\infty\),
\begin{equation}\label{eq:stein-to-Sd}
R_{d,\nu,\nu\lambda}(h)
\longrightarrow
R^{S,d}_{\sigma^2,\lambda}(h)
=
\frac{\sigma^2}{(1+\lambda|h|)^p},
\qquad h\in\mathbb R,
\end{equation}
whereas, as \(\nu\downarrow0\),
\begin{equation}\label{eq:stein-small-nu-limit}
R_{d,\nu,\nu\lambda}(h)
\longrightarrow
R^{S,2}_{\sigma^2,\lambda}(h)
=
\frac{\sigma^2}{1+\lambda|h|},
\qquad h\in\mathbb R.
\end{equation}
The convergence in \eqref{eq:stein-small-nu-limit} is uniform on every
compact subset of \(\mathbb R\). In particular, the small-\(\nu\)
closure is independent of \(d\). For convenience, write
\[
S_{\sigma^2,\lambda}
:=
S^{(2)}_{\sigma^2,\lambda}.
\]

\item[(iv)]
The small-\(\nu\) limiting process \(S_{\sigma^2,\lambda}\) exhibits
LRD and, for every \(r<v\) and \(s<t\),
\begin{equation}\label{eq:stein-small-nu-lrd}
\begin{aligned}
&\lim_{T\to\infty}T^3
\operatorname{Cov}\left(
S_{\sigma^2,\lambda}(v)
-
S_{\sigma^2,\lambda}(r),
\right.\\[-1mm]
&\hspace{4.3cm}\left.
S_{\sigma^2,\lambda}(t+T)
-
S_{\sigma^2,\lambda}(s+T)
\right)\\
&\qquad
=
-\frac{2\sigma^2}{\lambda}(v-r)(t-s).
\end{aligned}
\end{equation}

\item[(v)]
The scaling \(\rho=\nu\lambda\) is essential for the nondegenerate
small-\(\nu\) closure. If instead \(\rho>0\) is kept fixed, then
\begin{equation}\label{eq:stein-fixed-rho-small-nu}
R_{d,\nu,\rho}(h)
\longrightarrow
R_{\mathrm{nug}}(h)
:=
\begin{cases}
\sigma^2, & h=0,\\
0, & h\neq0,
\end{cases}
\qquad \nu\downarrow0.
\end{equation}
Thus, the fixed-\(\rho\) limit is the nugget covariance and does not
exhibit LRD.

\item[(vi)]
When \(d=2\), one has \(p=1\), and the reduction is exact for every
\(\nu>0\) and \(\rho>0\):
\begin{equation}\label{eq:stein-d2-exact}
R_{2,\nu,\rho}(h)
=
\frac{\sigma^2\nu}{\nu+\rho|h|}
=
\frac{\sigma^2}{1+(\rho/\nu)|h|}.
\end{equation}
Equivalently,
\[
R_{2,\nu,\rho}
=
R^{S,2}_{\sigma^2,\rho/\nu}.
\]
In particular, if \(\rho=\nu\lambda\), then
\[
R_{2,\nu,\nu\lambda}(h)
=
\frac{\sigma^2}{1+\lambda|h|}
\]
for every \(\nu>0\), so the small- and large-\(\nu\) closures coincide.

\item[(vii)]
When \(d=3\), one has \(p=3/2\), and
\begin{equation}\label{eq:stein-d3-covariance}
R_{3,\nu,\rho}(h)
=
\sigma^2
\frac{\Gamma\!\left(\nu+\frac32\right)}{\Gamma(\nu)}
\frac{\Gamma(\nu+\rho|h|)}
{\Gamma\!\left(\nu+\rho|h|+\frac32\right)}.
\end{equation}
Under \(\rho=\nu\lambda\), its two boundary closures are
\begin{equation}\label{eq:d3-small-nu}
R_{3,\nu,\nu\lambda}(h)
\longrightarrow
\frac{\sigma^2}{1+\lambda|h|},
\qquad \nu\downarrow0,
\end{equation}
and
\begin{equation}\label{eq:d3-large-nu}
R_{3,\nu,\nu\lambda}(h)
\longrightarrow
\frac{\sigma^2}{(1+\lambda|h|)^{3/2}},
\qquad \nu\to\infty.
\end{equation}
Thus, the closure of the \(d=3\) family contains both
\(S_{\sigma^2,\lambda}\) and \(S^{(3)}_{\sigma^2,\lambda}\), where
\begin{equation}\label{eq:S3-covariance}
R^{S,3}_{\sigma^2,\lambda}(h)
=
\frac{\sigma^2}{(1+\lambda|h|)^{3/2}}.
\end{equation}
Their covariance tails are proportional to \(|h|^{-1}\) and
\(|h|^{-3/2}\), respectively. Consequently, their LRD normalizations
are \(T^3\) and \(T^{7/2}\). In particular,
\begin{equation}\label{eq:S3-lrd-limit}
\begin{aligned}
&\lim_{T\to\infty}T^{7/2}
\operatorname{Cov}\left(
S^{(3)}_{\sigma^2,\lambda}(v)
-
S^{(3)}_{\sigma^2,\lambda}(r),
\right.\\[-1mm]
&\hspace{4.3cm}\left.
S^{(3)}_{\sigma^2,\lambda}(t+T)
-
S^{(3)}_{\sigma^2,\lambda}(s+T)
\right)\\
&\qquad
=
-\frac{15}{4}\sigma^2\lambda^{-3/2}
(v-r)(t-s).
\end{aligned}
\end{equation}
\end{enumerate}
\end{cor}

\begin{proof}
Since \(p=d/2>0\), the gamma-integral representation gives
\[
(1+\lambda|h|)^{-p}
=
\frac{1}{\Gamma(p)}
\int_0^\infty
u^{p-1}e^{-u}e^{-\lambda u|h|}
\,du.
\]
For every \(u>0\), the function
\[
h\longmapsto e^{-\lambda u|h|}
\]
is positive definite on \(\mathbb R\). Hence
\eqref{eq:Sd-covariance} is a positive mixture of covariance functions,
which proves~(i).

For \(h>0\),
\[
\frac{d^2}{dh^2}
R^{S,d}_{\sigma^2,\lambda}(h)
=
\sigma^2p(p+1)\lambda^2
(1+\lambda h)^{-p-2}
\sim
\sigma^2p(p+1)\lambda^{-p}h^{-p-2}.
\]
Applying the double-integral identity
\eqref{eq:increment-double-integral}, together with the uniformity of
this asymptotic relation under bounded shifts of its argument, yields
\eqref{eq:Sd-lrd-limit}. Therefore, the adopted LRD criterion holds
with
\[
b=p+2=\frac d2+2,
\]
proving~(ii).

For the large-\(\nu\) limit, substituting
\(\rho=\nu\lambda\) gives
\[
R_{d,\nu,\nu\lambda}(h)
=
\sigma^2
\frac{\Gamma(\nu+p)}{\Gamma(\nu)}
\frac{
\Gamma\!\left(\nu(1+\lambda|h|)\right)
}{
\Gamma\!\left(\nu(1+\lambda|h|)+p\right)
}.
\]
For fixed \(p>0\) and \(h\in\mathbb R\), the gamma-ratio asymptotics
yield
\[
\frac{\Gamma(\nu+p)}{\Gamma(\nu)}
\sim\nu^p
\]
and
\[
\frac{
\Gamma\!\left(\nu(1+\lambda|h|)\right)
}{
\Gamma\!\left(\nu(1+\lambda|h|)+p\right)
}
\sim
\nu^{-p}(1+\lambda|h|)^{-p}.
\]
Their product proves \eqref{eq:stein-to-Sd}.

For the small-\(\nu\) limit, let
\[
B(x,p)
:=
\frac{\Gamma(x)\Gamma(p)}{\Gamma(x+p)},
\qquad x,p>0.
\]
Then
\[
\frac{R_{d,\nu,\nu\lambda}(h)}{\sigma^2}
=
\frac{
B\!\left(\nu(1+\lambda|h|),p\right)
}{
B(\nu,p)
}.
\]
Define
\[
F_p(x):=xB(x,p).
\]
Since
\[
x\Gamma(x)\longrightarrow1
\qquad\text{and}\qquad
\Gamma(x+p)\longrightarrow\Gamma(p)
\]
as \(x\downarrow0\), one has
\[
F_p(x)\longrightarrow1.
\]
Writing \(c(h):=1+\lambda|h|\), we obtain
\[
\frac{
B\!\left(\nu c(h),p\right)
}{
B(\nu,p)
}
=
\frac{F_p\!\left(\nu c(h)\right)}{F_p(\nu)}
\frac{1}{c(h)}
\longrightarrow
\frac{1}{1+\lambda|h|},
\]
which proves \eqref{eq:stein-small-nu-limit}.

To verify local uniform convergence, fix \(H>0\). For \(|h|\leq H\),
\[
1\leq c(h)\leq1+\lambda H.
\]
Hence
\[
0<\nu c(h)\leq\nu(1+\lambda H),
\]
and \(F_p(x)\to1\) uniformly for
\(0<x\leq\nu(1+\lambda H)\) as \(\nu\downarrow0\). Therefore,
\eqref{eq:stein-small-nu-limit} is uniform on \([-H,H]\). Applying
\eqref{eq:Sd-lrd-limit} with \(d=2\), and hence \(p=1\), gives
\eqref{eq:stein-small-nu-lrd}, proving~(iii) and~(iv).

Suppose now that \(\rho>0\) remains fixed. For \(h\neq0\),
\[
\frac{\Gamma(\nu+p)}{\Gamma(\nu)}
\sim
\Gamma(p)\nu,
\qquad \nu\downarrow0,
\]
whereas
\[
\frac{\Gamma(\nu+\rho|h|)}
{\Gamma(\nu+\rho|h|+p)}
\longrightarrow
\frac{\Gamma(\rho|h|)}
{\Gamma(\rho|h|+p)}.
\]
Thus,
\[
R_{d,\nu,\rho}(h)\longrightarrow0,
\qquad h\neq0.
\]
At \(h=0\),
\[
R_{d,\nu,\rho}(0)=\sigma^2
\]
for every \(\nu>0\), proving
\eqref{eq:stein-fixed-rho-small-nu}. For the resulting nugget
covariance, the covariance between the two distant increments is zero
for all sufficiently large \(T\); hence it does not satisfy the adopted
LRD criterion. This proves~(v).

If \(d=2\), then \(p=1\), and the recurrence relation
\[
\Gamma(x+1)=x\Gamma(x)
\]
applied to \(R_{d,\nu,\rho}\) gives
\eqref{eq:stein-d2-exact}, proving~(vi).

Finally, if \(d=3\), then
\[
p=\frac32,
\qquad
p+2=\frac72,
\qquad
p(p+1)
=
\frac32\frac52
=
\frac{15}{4}.
\]
Substitution into the general Stein covariance gives
\eqref{eq:stein-d3-covariance}. Equations
\eqref{eq:d3-small-nu} and \eqref{eq:d3-large-nu} follow from
\eqref{eq:stein-small-nu-limit} and \eqref{eq:stein-to-Sd},
respectively. Finally, substituting \(p=3/2\) into
\eqref{eq:Sd-covariance} and \eqref{eq:Sd-lrd-limit} yields
\eqref{eq:S3-covariance} and \eqref{eq:S3-lrd-limit}.
\end{proof}

\end{document}
